% !TeX root = ../main_en.tex
\subsection{Motivation}
\label{subsec:motivacion-contexto}

The autonomous trajectory tracking of a quadcopter requires having a sufficient estimate of its state with respect to the environment and the reference, and transforming the tracking error into an actuation that stabilizes the vehicle while guiding it along the designated path. The difficulty does not lie in minimizing an instantaneous error with respect to a fixed position; rather, the reference evolves over time, and the controller acts on a rigid body with coupled dynamics and exposure to internal and external disturbances. Therefore, it must coordinate translation and attitude without losing stability or failing in the displacement.

This coupled nature motivates the use of a two-level control architecture \cite{beard2012small,mellinger2011minimum,lee2010geometric}. In this case, a cascaded PD first transforms position and velocity errors into a desired force in the local East--North--Up (ENU) reference frame. Subsequently, the inner loop converts this force, along with the yaw reference, into a desired attitude, thrust, and moments in the Forward--Right--Down (FRD) body reference frame. Finally, the mixer distributes these commands among the four rotors. Tuning a single controller for maneuvers of a diverse nature can impose a compromise between, for example, maintaining a position and tracking a geometry with continuously changing motion. Specialization avoids this compromise but scales the number of controllers to be tuned and validated, as well as the constant switching of gains during operation. For this reason, in this work, it is decided to build a set of controllers with a PD for each trajectory family, and it is studied whether a neural implementation can condense the joint operation of the experts. This idea of having different control parameters depending on the operating regime is closely related to gain scheduling \cite{leith2000survey}.

The comparison between the two approaches requires an environment where the model, references, disturbances, random seeds, actuation limits, and metrics remain under control. The developed 6DOF simulator provides this testbed where all controllers share the vehicle, state observation, and actuation chain. Furthermore, each run preserves traceability of the configuration to reproduce the experiment, as well as the telemetry necessary for subsequent review. This platform does not intend to replace a future in-flight validation, but rather to isolate the differences in a controlled and reproducible environment.

\begin{figure}[H]
  \centering
  \includegraphics[width=\textwidth]{FIG-001}
  \caption{Conceptual diagram of the experimental workflow of the work. It consists of trajectory generation, specialization of classic controllers, training of neural networks, and comparative evaluation. Source: own elaboration.}
  \label{fig:estrategia-experimental}
\end{figure}

In Figure \ref{fig:estrategia-experimental}, the experimental logic structured in two stages is detailed visually. The first stage illustrates obtaining the expert controllers and generating data to train the neural networks. The second stage shows the comparative evaluation of the three control references in the simulator under the same conditions: conditions within the training distribution or \textit{in-distribution} (ID), and conditions outside said distribution or \textit{out-of-distribution} (OOD).

\subsubsection*{Problem Addressed}

The main problem addressed in this work consists of determining whether a single controller using neural networks, trained on data generated by PD controllers optimized for different trajectory families, can replace this set of controllers when evaluated under the same references and disturbances. The comparison should indicate whether the neural implementation maintains competitive behavior under known conditions and whether it possesses better transfer to composite, new, and varied trajectories.

Three references are compared. The \emph{specialized PD} is the controller whose parameters have been fixed for the family to which the trajectory belongs. A \emph{transferred PD} is one of those same controllers applied, without retuning, to another family. The neural system replaces the outer loop and performs a \emph{desired force prediction} from the observation and the reference, producing a three-dimensional vector in the ENU reference frame and preserving the classic inner loop. In this way, the network does not transmit commands directly to the rotors nor does it generate moments in the FRD reference frame.

On one hand, \emph{competitive performance} is defined as the ability to preserve, in closed-loop, the completion of the trajectory while obtaining tracking errors, control effort, and protection usage comparable to those of the relevant reference. The metrics and evaluation criteria associated with these concepts are defined in the \hyperref[subsubsec:metricas-criterios]{epigraph on metrics, success criteria, and analysis}. On the other hand, \emph{transfer} is evaluated by applying the controller outside the operating conditions for which it was originally tuned. To structure this analysis, the evaluation distinguishes between in-distribution (ID, \textit{in-distribution}) and out-of-distribution (OOD, \textit{out-of-distribution}) scenarios, whose attributes are defined later in \apartadoref{subsec:flujo-reproducibilidad}. In the ID case, scenarios of the same type as those used to tune the PDs and train the networks are evaluated, albeit with parameters reserved for testing. In the OOD case, unseen dynamic variations, external disturbances, or new geometries are imposed. Within the latter geometric references, a distinction is made between compositions of already known trajectories and completely unseen geometric families.

It is important to point out that this study does not intend to demonstrate the complete superiority of any particular architecture, nor to guarantee the robustness of the system against any physical disturbance or its suitability for real flight. Likewise, the neural controller is not credited with the attitude stabilization of the inner loop nor the management of protections imposed through saturations and force limits. The scope of this work is limited to a traceable experimental comparison, carried out on the same physical model and environment, between gain specialization, classic controller transfer, and neural imitation.

\subsection{Hypotheses and Objectives}
\label{subsec:hipotesis-objetivos}

The central hypothesis holds that a single neural system, trained by imitation to predict the desired force of the outer loop, can achieve competitive tracking capacity compared to the specialized PD of each family under known conditions and transfer better than the PDs applied outside their native family, while preserving the classic inner loop in both.

The general objective is to develop a reproducible simulation environment that allows determining whether the neural imitation controller can replace a set of specialized PD controllers and offer better performance in trajectories not used during gain tuning and training.

This objective is broken down into the following specific objectives:

\begin{enumerate}
  \item Formulate and implement a 6DOF rigid-body model for a quadcopter, with appropriate reference frames, actuators, simplified disturbances, and differentiated update frequencies;
  \item Define trajectory families and parameters that allow repeating the experimental conditions;
  \item Implement a cascaded PD and obtain, through a deterministic tuning process over the training scenarios, a specialized controller for each family;
  \item Assess the specialization of the PDs and their transfer capability by applying each controller to families other than its native one;
  \item Train MLP, GRU, and LSTM to predict the desired force and compare the absence or presence of temporal memory under the same conditions;
  \item Evaluate the quality of the prediction in supervised execution and, separately, the closed-loop tracking in ID and OOD conditions; and
  \item Preserve sufficient results to reconstruct the experiments.
\end{enumerate}

\subsection{Structure of the Document}
\label{subsec:estructura-documento}

The contribution of the work is the integration of four elements: a 6DOF simulator, a set of trajectories and parameters through scenarios, a cascaded PD selection process specialized by family, and a closed-loop comparison of MLP, GRU, and LSTM. The interest lies in building an environment that allows distinguishing which properties originate from the specialized PD, what the network learns, what the classic inner loop contributes, and what the limits of this approach to the problem are.

After this introduction, the document is structured in four blocks: the state of the art on modeling, classic control, and imitation control; the methodology, which details the simulator and the evaluated control architectures; the analysis of experimental results; and, finally, the conclusions and future lines of development.
